
\documentclass[a4paper,11pt]{letter} 
%%%% choose one of the following 4 lines to set colour and/or  date style, 
\usepackage[colour]{leedsletter} % Day, num Month year, and colour (use lp2)
%\usepackage[]{leedsletter} % Black-and-white, with US style date Month num, year
%\usepackage[myukdate]{leedsletter} 
%\usepackage[colour]{leedsletter} 
%%%% end
\department{ \vspace{-1cm} {\bf Department of Political Science}\\[38pt]
\vspace{-1.5cm}
~\\
Pennsylvania State University\\
203 Pond Lab \\
University Park, PA 16802\\ \vspace{.1cm} \\
%%%% edit the following three lines to suit you.  If you add a line, then 38pt -> 26pt above
T~~ 814.865.4598\\
F~~ 814.863.8979\\
E~~ bdesmarais@psu.edu  \\
W~~brucedesmarais.com\\ \vspace{-2cm}} 
%%%% end

\alignit{10pt} %%  tweaking parameter - probably no need to change. 
\normalsize

%%%  need to edit  user info 
\signature{Bruce A. Desmarais Ph.D.\\[5pt] 
\rlap{DeGrandis-McCourtney Early Career Professor} \\
\rlap{Director, Graduate Programs in Social Data Analytics} \\
\rlap{Director, Center for Social Data Analytics} \\ Department of Political Science\\ \rlap{Institute for Computational and Data Sciences}}
% end personal details 

% end personal details 
%%%%%%%%%%%%%%%%%%%%%% a bove this line is sender specific %%%%%%%%%%%%
%%%%%%%%%%%%%%%%%%%%%% and probably only needs edited once %%%%%%%%%%%%
%%% what follows is letter specific


% the next line is specific to this letter, and may be used for "short" letters 

\skipit{.6}  % this amount (default = 1.7) increases 2 gaps: address-date-opening 
 \vspace{-3.5cm}
\begin{document} 
\begin{letter} 


%% if this is a ``short'' or long address, then alter \skipit above 
 

 
\opening{Dear search committee members,}
\medskip 

%% Todo
% revise first two paragraphs

I write to express great interest in the Canada Excellence Research Chair in Computational Social Science, AI and Democracy at McGill University. The driving passion in my career is to push boundaries regarding the data and methodologies that researchers have at their disposal to aid in understanding  society. To that end, I develop and apply cutting edge methods of machine learning and network science; teach courses on advanced methods of machine learning and data science; and build interdisciplinary institutions to support the emerging fields of data science and computational social science. It is clear that artificial intelligence presents the scientific community and society with unprecedented opportunities, risks, and ethical considerations. It is inspiring to see McGill investing resources to position itself at the forefront of AI-related research and training, especially with an emphasis on Democracy. I focus on big-picture projects, in which I seek external funding to support large teams of statistical, social, computational, and information scientists to gather unique big data resources, and develop and apply new statistical and computational methods that allow us to advance understanding of the political world. I have been fortunate to work on the founding team for three degree programs and two research centers situated at the intersection of social and computational sciences. I would welcome opportunities to participate in and support interdisciplinary institutions related to computational social science, AI, and democracy at McGill.

In my research I develop and apply advanced statistical and computational methods with an emphasis on substantive problems in political science and public policy, such as policy innovation and diffusion, and digital engagement with political elites. In my most recent and ongoing research, I am focused on the study of online engagement with and by governments and public officials. I am an early adopter of computational tools that form the basis of generative AI technology. For example, we published a paper in 2021 in which we used a large language model to identify Tweets about the COVID-19 pandemic that were posted by subnational elected officials in the U.S. My recent substantive research addresses the intersection of digital politics and public policy, and we consistently draw upon generative AI tools for applications of computer vision, natural language processing, and network modeling. I study online social networks (e.g, X/Twitter, Facebook, YouTube) as modern public forums in which policymakers engage with each other, and other stakeholders. We study problems such as understanding the connections among individual policymakers online, and the dissemination of misinformation by policymakers themselves on X/Twitter and Facebook. This focus on policymakers' online behavior was motivated in large part by my long-term research agenda on policy networks. My collaborators and I have focused on (1) measuring policy innovation, and (2) modeling the spread of policies within the network science framework. We have provided new data sources on state policies, and software for measuring policies and modeling policy diffusion networks.  The interdisciplinary nature of my research fits this position well, and I would naturally engage with interdisciplinary efforts focused on CSS, AI, and democracy. I have published in leading outlets across several disciplines, including political science, public policy, statistics, sociology, computer science, neuroscience, and physics; and have led highly interdisciplinary project teams. 


As a teacher, I am experienced in both undergraduate and graduate courses relevant to this role. Most of my teaching has been in quantitative research methods. I have taught undergraduate courses in {\em Quantitative Methods} and {\em Research Design for Social Data Analytics}. I have taught graduate introductory statistical methods courses, and several advanced courses, including {\em Machine Learning}, {\em Network Analysis}, {\em  Approaches and Issues in Social Data Analytics}, and {\em Big Social Data}.  My graduate machine learning course includes a four-week module on deep learning and generative AI.  My approach to research methods instruction reflects the movement for, ``Data Science for Everyone.'' This is a philosophy of data science instruction that establishes a top priority of making the coursework accessible to all who are interested in learning, regardless of their math or computing background, and with careful attention to how individuals' identities may affect their experience of or perspective on the course content. With each new semester I revise my syllabus in ways that lower barriers to entry, while maintaining high standards of rigor. 

I am excited to see the emphasis in this position on the leadership of a lab. Given the interdisciplinary and less established nature of career paths in computational social science, strong mentorship and collaboration are crucial to the success of junior researchers. I conduct my research according to a lab model, in which I work actively with undergraduate students, graduate students, research faculty, and postdocs; and I take responsibility for funding their work. We hold weekly lab meetings to foster connections, review recent advancements in the literature, and provide ample opportunities for lab members to seek feedback. I have advised nine PhD students who have completed their PhDs, and have placed in tenure track roles at research universities, interdisciplinary data science postdocs, and senior data scientist positions at Meta and Amazon. I have supervised four postdocs, two of whom are currently working with me in positions at Penn State, and two who have gone on to roles in academia and industry. I am impressed with the research centers in politics, public policy, and computer science at McGill that are connected with this opportunity, and would be excited for my lab to serve as a point of integration among these institutions. 

I am highly motivated in applying to this opportunity at McGill, as it uniquely brings together the three central pillars of my work—computational methods, public policy, and democratic institutions—with significant resources and a clear mandate for interdisciplinary leadership. I would be excited to help shape and sustain research infrastructure that bridges centers and groups at McGill, and to drive the mentorship and collaboration that strengthen the broader research community around the Chair. Personally, the move to Montreal is an exciting prospect for my family as well; with ancestral roots in the region and strong family connections to the Northeast U.S., we would welcome the opportunity to move to the Montreal metro area. I would be honored to bring my experience and vision to McGill to help advance its ambitious goals at the intersection of CSS, AI and democracy.

\closing{Sincerely,} 
\end{letter} 
\end{document} 
